\clearpage
\VerdierPhysicalPageAnchor{0161}%
\begingroup\small\setlength{\parskip}{0pt}\setlength{\emergencystretch}{1.5em}
\VerdierSemanticAnchor{ch3-3-2-9-exactness-test}
is exact in $\mathcal D^\wedge$. The sequence (3.2.9.1) is exact if and only
if, for every object $W$ in $\mathcal D$, the sequence
\[
\operatorname{Hom}_{\mathcal D^\wedge}(h(W),h(X))\longrightarrow
\operatorname{Hom}_{\mathcal D^\wedge}(h(W),h(Y))\longrightarrow
\operatorname{Hom}_{\mathcal D^\wedge}(h(W),h(Z))
\]
is exact. This follows from full faithfulness of $h$ and (1.2.1).

\VerdierSemanticAnchor{ch3-3-2-9-factorization}
Let $\mathcal A$ be an abelian category, and let
\[
\mathcal E x(\mathrm A(\mathcal D),\mathcal A)
\longrightarrow\mathcal H\!om(\mathcal D,\mathcal A)
\]
be the functor “composition with $\mathcal H$.” This functor is clearly fully
faithful: a morphism between exact functors is uniquely determined once its
values on sufficiently many projective objects are known. It remains to show
that every cohomological functor
$\Psi\colon\mathcal D\to\mathcal A$ factors through
$\mathrm A(\mathcal D)$:
\VerdierSemanticAnchor{ch3-3-2-9-2}
\[
\begin{tikzcd}[column sep=5em,row sep=3em]
\mathcal D\arrow[r,"\mathcal H"]\arrow[dr,"\Psi"']&
\mathrm A(\mathcal D)\arrow[d,"\Phi"]\\
&\mathcal A
\end{tikzcd}
\tag{3.2.9.2}
\]
with $\Phi$ exact. For every object
$X\xrightarrow fY$ of $\mathrm A(\mathcal D)$, set
\begingroup\renewcommand{\thefootnote}{(*)}%
\[
\Phi(X\xrightarrow fY)=
\operatorname{Image}_{\mathcal A}
(\Psi(X)\xrightarrow{\Psi(f)}\Psi(Y)).\footnote{Note that the construction
of $\Phi$ involves choosing an “image” functor in $\mathcal A$. This choice
is, of course, unique up to unique isomorphism. We may always require, as we
shall henceforth, that the “image” functor associate to every identity
morphism $\operatorname{id}_T$ in $\mathcal A$ the object $T$ equipped with
the morphism $T\xrightarrow{\operatorname{id}_T}T$.}
\]
\endgroup
This defines a functor from $\mathrm A(\mathcal D)$ to $\mathcal A$ for
which diagram (3.2.9.2) is commutative. We show that $\Phi$ is exact. It
follows immediately from (3.2.7) that $\Phi$ is left exact. Passing to
opposite categories shows in the same way that $\Phi$ is right exact. This
completes the proof of Theorem (3.2.1).

\VerdierSemanticAnchor{ch3-3-2-10}
\medskip\noindent\textbf{Remark 3.2.10.}\quad
1) The universal property described in Theorem (3.2.1, c)) determines
$\mathrm A(\mathcal D)$ up to equivalence.

2) When countable direct sums or countable products are representable in
$\mathcal D$, all projective (respectively, injective) objects of
$\mathrm A(\mathcal D)$ are isomorphic to objects of $\mathcal D$. Indeed,
such an object is a direct summand in $\mathrm A(\mathcal D)$ of an object
of $\mathcal D$. It is therefore the kernel of a projector of an object of
$\mathcal D$. Under the stated hypotheses, projectors are decomposable in
$\mathcal D$ (1.2.9).
\endgroup
